\documentclass[11pt,a4paper]{article}
\usepackage[T1]{fontenc}
\usepackage{lmodern,amsmath,amssymb,amsthm,mathtools}
\usepackage[margin=27mm]{geometry}
\usepackage[colorlinks=true,linkcolor=blue,urlcolor=blue]{hyperref}
\newtheorem{theorem}{Theorem}
\newtheorem{lemma}[theorem]{Lemma}
\newtheorem{corollary}[theorem]{Corollary}
\newcommand{\Q}{\mathbb Q}\newcommand{\Z}{\mathbb Z}
\title{Square abscissas and direct charts for the elliptic subgroup}
\author{}\date{8 November 2026}
\begin{document}\maketitle
\begin{abstract}
Under the stated elliptic group hypothesis, two square equations
describe exactly the nonzero points of the cyclic subgroup used
in the arithmetic construction. No parity condition on the index
is necessary. A cubic change of coordinates gives one rational
quartic presentation that covers every point, with no choice among
component charts.
\end{abstract}

\section{The subgroup described by a square abscissa}

Put
\[
 K=\Q(\sqrt2),\qquad d=3-\sqrt2,\qquad D=d^2=11-6\sqrt2,
 \qquad E:y^2=x^3-Dx.
\]
Let $O$ denote the identity of $E$, and set
\[
 R_0=(1+\sqrt2,2\sqrt2),\qquad P=2R_0.
\]
The arithmetic hypothesis in this note is
\begin{equation}\label{eq:group}
                 2E(K)=\Z P.
\end{equation}
We also assume, as in the source construction, that $P$ has infinite
order. Thus every $Q\in E(K)$ has a presentation
$Q=mR_0+T$ with $m\in\Z$ and $T\in E[2](K)$.

\begin{lemma}\label{lem:descent}
The map
\[
 \delta:E(K)\longrightarrow K^\times/K^{\times2},\qquad
 \delta(O)=1,\quad \delta((0,0))=-D,\quad
 \delta(Q)=x(Q)\quad(x(Q)\ne0),
\]
is a group homomorphism.
\end{lemma}
\begin{proof}
Consider a nonvertical line $y=\lambda x+\mu$. Its intersections
with $E$, counted with multiplicity, have abscissas given by
\[
 x^3-\lambda^2x^2-(D+2\lambda\mu)x-\mu^2=0.
\]
When $\mu\ne0$, the product of the three abscissas is $\mu^2$.
When $\mu=0$, one intersection is $(0,0)$ and the product of the
other two abscissas is $-D$. In either case the product of the
three $\delta$-classes is trivial. Vertical lines meet $Q$,
$-Q$, and $O$, whose classes also have trivial product. These
identities, including tangent lines through their intersection
multiplicities, give the assertion by the chord and tangent law.
\end{proof}

\begin{theorem}\label{thm:chart}
Assume \eqref{eq:group}. The points arising from
\begin{equation}\label{eq:chart}
             x=h^2,\qquad j^2=h^4-D,\qquad y=hj,
             \qquad h,j\in K,
\end{equation}
are exactly $E(K)\cap(\Z P\setminus\{O\})$.
Every such point has $h\ne0$ and $j\ne0$.
\end{theorem}
\begin{proof}
Both real embeddings of $d$ and $D$ are positive. Consequently
$h=0$ would imply $j^2=-D$, which is impossible. If $j=0$,
then $h^2=d$ or $h^2=-d$. The second alternative is excluded
by signs. The first is excluded by $N_{K/\Q}(d)=7$, since the
norm of a square in $K$ is a rational square. Thus the point
$Q=(h^2,hj)$ is affine and is not a point of order two.

Write $Q=mR_0+T$ as above. The four two-torsion points and their
classes are
\[
\begin{array}{c|cccc}
T&O&(0,0)&(d,0)&(-d,0)\\ \hline
\delta(T)&1&-D&d&-d.
\end{array}
\]
Every representative in the second row has the same sign at the
two real embeddings. In contrast,
$\delta(R_0)=1+\sqrt2$ has opposite signs at those embeddings.
Since $\delta(Q)$ is a square class, the homomorphism property
forces $m$ to be even. Write $m=2k$. Then
$Q=kP+T$, and $\delta(P)=\delta(2R_0)=1$ implies
$\delta(Q)=\delta(T)$. Of the four listed classes only $1$ is
a square: $-D$ and $-d$ are negative, while $d$ has nonsquare
norm. Hence $T=O$ and $Q=kP$, where $k\ne0$ because $Q$ is
affine.

Conversely, if $Q=kP$ with $k\ne0$, then $Q=2(kR_0)$ and
$\delta(Q)=1$. The point is nontorsion, so $x(Q)\ne0$. Choose
$h\in K^\times$ with $h^2=x(Q)$ and put $j=y(Q)/h$.
The equation of $E$ gives $j^2=h^4-D$, as required.
\end{proof}

\section{A rational quartic presentation}

Write $h=h_0+h_1\sqrt2$ and
\[
 U=h_0^4+12h_0^2h_1^2+4h_1^4-11,
 \qquad V=4h_0h_1(h_0^2+2h_1^2)+6.
\]
Thus $h^4-D=U+V\sqrt2$. Define
\begin{align*}
 a&=2U+V, & b&=-(U+V),\\
 A&=6U^3-9UV^2-4V^3,
 &B&=-4U^3+6UV^2+3V^3.
\end{align*}
Then the square-root equation for $j$ is equivalent to
\begin{equation}\label{eq:quartic}
              Y^4-4AY^2+8B^2=0,\qquad Y\ne0.
\end{equation}
Every such rational root recovers $j$ by
\begin{equation}\label{eq:recovery}
 j=\frac{a-b\sqrt2}{2U^2-V^2}
                   \left(\frac Y2+\frac BY\sqrt2\right).
\end{equation}

\begin{corollary}\label{cor:single-chart}
Under \eqref{eq:group}, the formulas above represent all nonzero
points of $\Z P$ by two rational base coordinates $h_0,h_1$
and one finite quartic root $Y$. Every permitted rational root
is simple. No finite choice of chart is required.
\end{corollary}
\begin{proof}
Theorem~\ref{thm:chart} reduces the assertion to the square-root
presentation. If $j=u+v\sqrt2\ne0$, then
$U=u^2+2v^2$ and $V=2uv$, and direct expansion gives
\[
 (a+b\sqrt2)j
       =2(u^3-2v^3)+(4v^3-u^3)\sqrt2.
\]
The rational component is nonzero because $2$ is not a rational
cube. The multiplier has norm $2U^2-V^2\ne0$, since $(U,V)\ne(0,0)$
and $2$ is not a rational square. Its square times $j^2$ is
$A+B\sqrt2$. If $Y$ is twice the rational component of
$(a+b\sqrt2)j$, these identities give \eqref{eq:quartic} and
\eqref{eq:recovery}.

Conversely, a nonzero root $Y$ gives
\[
 \left(\frac Y2+\frac BY\sqrt2\right)^2=A+B\sqrt2
                     =(a+b\sqrt2)^2(U+V\sqrt2),
\]
so \eqref{eq:recovery} gives the required square root $j$.
At such a root the derivative of the quartic is
\[
 4Y^3-8AY=8Y\,N_{K/\Q}\left(\frac Y2+\frac BY\sqrt2\right)\ne0.
\]
This proves the simple-root assertion as well.
\end{proof}

The condition $Y\ne0$ must be retained when this presentation
is used: the quartic equation alone does not justify the displayed
rational recovery formula at $Y=0$. In a finite-algebra
construction it is imposed by the localization step, with its
unit condition checked separately. The general version of this
square-root chart, valid for every quadratic extension of a number
field, is proved in the companion paper \emph{A single quartic
chart for square roots in a quadratic extension}.

\end{document}
